\documentclass[10pt]{article}
\usepackage[T1]{fontenc}
\usepackage[utf8]{inputenc}
\usepackage[french]{babel}
\usepackage{amsmath,amssymb,mathrsfs,mathpazo}
\usepackage{geometry}
\geometry{a4paper,margin=2cm}
\usepackage{graphicx,array,tabularx,tabularray,colortbl,multirow,multicol,booktabs,enumitem,pifont,ulem,fancybox,diagbox,cancel}
\usepackage{pstricks,pstricks-add,pst-plot,pst-tree,pst-node,pst-text,pst-func,pst-math,pst-bspline,pst-3dplot,pst-eucl,pst-solides3d}
\let\pscalculate\undefined
\usepackage{pst-bezier}
\usepackage{tikz,pgfplots,tkz-tab}
\usetikzlibrary{arrows,arrows.meta,patterns,calc,positioning,trees,decorations.markings,shapes.geometric,shapes.misc,shapes.arrows,intersections,decorations.pathmorphing,decorations.pathreplacing,backgrounds,shadows,babel}
\pgfplotsset{compat=1.18}
\usepgfplotslibrary{fillbetween}
\usepackage[np]{numprint}
\usepackage{listings,makecell,pas-tableur,twoopt}
\let\barre\undefined
\usepackage[pstricks]{tablvar}
\definecolor{darkgreen}{RGB}{0,100,0}
\providecommand{\red}{\color{red}}
\providecommand{\blue}{\color{blue}}
\providecommand{\green}{\color{green}}
\providecommand{\black}{\color{black}}
\providecommand{\white}{\color{white}}
\providecommand{\gray}{\color{gray}}
\providecommand{\cyan}{\color{cyan}}
\providecommand{\magenta}{\color{magenta}}
\providecommand{\yellow}{\color{yellow}}
\providecommand{\violet}{\color{violet}}
\providecommand{\orange}{\color{orange}}
\providecommand{\euro}{\texteuro}
\providecommand{\EUR}{\texteuro}
\providecommand{\ieme}{\textsuperscript{e}}
\providecommand{\degres}{\ensuremath{^\circ}}
\providecommand{\degree}{\ensuremath{^\circ}}
\providecommand{\vv}[1]{\overrightarrow{#1}}
\providecommand{\Coord}[1]{#1}
\providecommand{\decosix}{}
\providecommand{\starredbullet}{\ensuremath{\bullet}}
\pagestyle{empty}
\setlength{\parindent}{0pt}
\usepackage{adjustbox}\DeclareUnicodeCharacter{2236}{:}\DeclareUnicodeCharacter{2190}{\ensuremath{\leftarrow}}
\providecommand{\R}{\mathbb{R}}
\providecommand{\N}{\mathbb{N}}
\providecommand{\Z}{\mathbb{Z}}
\providecommand{\Q}{\mathbb{Q}}
\providecommand{\D}{\mathbb{D}}
\providecommand{\C}{\mathbb{C}}
\providecommand{\e}{{\mathrm{e}}}
\providecommand{\vect}[1]{\overrightarrow{#1}}
\providecommand{\vectt}[1]{\overrightarrow{\mathrm{#1}}}
\providecommand{\vvt}[1]{\overrightarrow{\mathrm{#1}}}
\providecommand{\np}[2][]{{#2~#1}}
\providecommand{\vv}[1]{{\overrightarrow{#1}}}
\renewcommand{\d}{{\mathrm{d}}}
\providecommand{\degres}{{^\circ}}
\providecommand{\degre}{{^\circ}}
\providecommand{\blue}{\color{blue}}
\providecommand{\red}{\color{red}}
\providecommand{\enskip}{\quad}
\providecommand{\llbracket}{\lbrack\!\lbrack}
\providecommand{\rrbracket}{\rbrack\!\rbrack}
\providecommand{\euro}{{\text{€}}}
\providecommand{\no}{{\text{n°}}}
\providecommand{\interval}[2]{{[#1;#2]}}
\providecommand{\intervFO}[2]{{[#1;#2[}}
\providecommand{\dotfill}{\ldots}

\providecommand{\ds}{\displaystyle}
\providecommand{\textding}[1]{#1}
\providecommand{\Oij}{\(\mathrm{O};\vec{\imath},\vec{\jmath}\)}
\providecommand{\Oijk}{\(\mathrm{O};\vec{\imath},\vec{\jmath},\vec{k}\)}
\providecommand{\Ouv}{\(\mathrm{O};\vec{u},\vec{v}\)}
\providecommand{\intervFF}[2]{\ensuremath{[#1;#2]}}
\renewcommand{\euro}{\ifmmode\text{\texteuro}\else\texteuro\fi}
\DeclareUnicodeCharacter{00B0}{\ifmmode^{\circ}\else\textdegree\fi}
\DeclareUnicodeCharacter{20AC}{\ifmmode\text{\texteuro}\else\texteuro\fi}
% Énoncés : Éducation nationale. Transcription et corrigés : APMEP, Denis Vergès et contributeurs. Découpage et classement : CoopMaths, Sébastien Lozano, Éric Elter et contributeurs. Adaptation de la présentation : Maths Explorer.
% CC BY-SA : https://coopmaths.fr/www/about/
% https://coopmaths.fr/alea/static/bac/2025/tex/bac_2025_11_sujet1_ameriquesud_1.tex
\begin{document}
\noindent\textit{APMEP / CoopMaths. CC BY-SA. Présentation adaptée par Maths Explorer.}\par\medskip
\medskip

Un étudiant mange tous les jours au restaurant universitaire. Ce restaurant propose des plats végétariens et des plats non végétariens.

\begin{itemize}[label=$\bullet~~$]
\item Lorsqu'un jour donné l'étudiant a choisi un plat végétarien, la probabilité qu'il choisisse un plat végétarien le lendemain est $0,9$.
\item Lorsqu'un jour donné l'étudiant a choisi un plat non végétarien, la probabilité qu'il choisisse un plat végétarien le lendemain est $0,7$.
\end{itemize}

\medskip

Pour tout entier naturel $n$, on note $V_n$ l'évènement \og l'étudiant a choisi un plat végétarien le $n^{\text{e}}$ jour \fg{} et $p_n$ la probabilité de $V_n$.

Le jour de la rentrée, l'étudiant a choisi le plat végétarien. On a donc $p_1 = 1$.

\medskip

\begin{enumerate}
\item
	\begin{enumerate}
		\item Indiquer la valeur de $p_2$.
		\item Montrer que $p_3 = 0,88$. On pourra s'aider d'un arbre pondéré.
		\item Sachant que le $3\up{\text{e}}$ jour l'étudiant a choisi un plat végétarien, quelle est la probabilité qu’il ait choisi un plat non végétarien le jour précédent ?

On arrondira le résultat à $10^{-2}$.
	\end{enumerate}
\item Recopier et compléter l'arbre pondéré ci-dessous:

\begin{center}
\pstree[treemode=R,levelsep=3cm]{\TR{}}
{\pstree{\TR{$V_n~$}\taput{$p_n$}}
	{\TR{$V_{n+1}$}
	\TR{$\overline{V_{n+1}}$} 
	}
\pstree{\TR{$\overline{V_n}~$}}
	{\TR{$V_{n+1}$}
	\TR{$\overline{V_{n+1}}$} 
	}
}
\end{center}
\item Justifier que, pour tout entier naturel $n \geqslant 1,\:\: p_{n+1} = 0,2p_n + 0,7$.
\item On souhaite disposer de la liste des premiers termes de la suite $(p_n)$ pour $n \geqslant 1$.

Pour cela, on utilise une fonction appelée \texttt{repas} programmée en langage Python dont on propose trois versions, indiquées ci-dessous.
\end{enumerate}
\hspace{-0.5cm}\vspace{-0.5cm}

\begin{center}
{\footnotesize
\begin{ttfamily}
\begin{tabularx}{\linewidth}{*{3}{>{\centering \arraybackslash}X}}
\multicolumn{1}{c}{Programme 1}&\multicolumn{1}{c}{Programme 2}&\multicolumn{1}{c}{Programme 3}\\ 
\begin{tabular}{l| l|}\cline{2-2}
1&def repas(n):\\
2&~~p=1\\
3&~~L=[p]\\
4&~~for k in range(1,n):\\
5&~~~~p = 0.2*p+0.7\\
6&~~~~L. append(p)\\
7&~~~return(L)\\ \cline{2-2}
\end{tabular}&
\begin{tabular}{l| l|}\cline{2-2}
1&def repas(n):\\
2&~~p=1\\
3&~~L=[p]\\
4&~~for k in range(1,n+1):\\
5&~~~~p = 0.2*p+0.7\\
6&~~~~L. append(p)\\
7&~~return(L)\\ \cline{2-2}
\end{tabular}&
\begin{tabular}{l| l|}\cline{2-2}
1&def repas(n):\\
2&~~p=1\\
3&~~L=[p]\\
4&~~for k in range(1,n):\\
5&~~~~p = 0.2*p+0.7\\
6&~~~~L.append(p+1)\\
7&~~return(L)\\ \cline{2-2}
\end{tabular}\\
\end{tabularx}
\end{ttfamily}
}
\end{center}

\begin{enumerate}[start=4]
	\item
	\begin{enumerate}
		\item Lequel de ces programmes permet d’afficher les $n$ premiers termes de la suite $(p_n)$ ? Aucune justification n'est attendue.
		\item Avec le programme choisi à la question \textbf{a.} donner le résultat affiché pour $n = 5$.
	\end{enumerate}
\item Démontrer par récurrence que, pour tout naturel $n \geqslant 1,\:\: p_n = 0,125 \times 0,2^{n - 1} + 0,875$.
\item En déduire la limite de la suite $(p_n)$.
\end{enumerate}

\bigskip

\end{document}
